\section{บทนำ}
อักษรไทยมีสระและวรรณยุกต์ซ้อนอยู่บนและล่างแนวพยัญชนะ เราเรียกอักขระซ้อนเหล่านี้ว่า mark คือสระบน สระล่าง และวรรณยุกต์ (รายการอยู่ในหัวข้อ~\ref{sec:metric}) mark ตัวเดียวเปลี่ยนความหมายและเสียงของคำได้ เช่น ``ข่าว'' กับ ``ข้าว'' ถ้าการถอดข้อความทำ mark หายหรือสลับ ผู้อ่านจะเห็นว่าผิด แม้อักขระอื่นเกือบถูกหมด

แต่ OCR ไทยมักวัดด้วยตัวชี้วัดของทั้งข้อความ เช่น ความทับซ้อนของ n-gram และ normalized edit distance ใน ThaiOCRBench \cite{nonesung2025thaiocrbench} หรือ character error rate ในงานอื่น ตัวชี้วัดเหล่านี้ไม่บอกว่าปัญหาอยู่ที่ mark หรือไม่ คะแนนระดับลำดับยังลงโทษโมเดลที่อ่านคอลัมน์ในลำดับต่างจากเฉลยด้วย

ปัจจุบันมี vision-language model (VLM) แบบเปิดสำหรับ OCR ไทยแล้ว \cite{nonesung2026typhoonocr,nectec2026pathumma,pipatanakul2026synthetic} เมื่อโมเดลทำ mark ผิด อธิบายได้สามทาง แต่ละทางมีวิธีแก้ต่างกัน
\begin{itemize}
\item หลักฐานในภาพของ mark เล็กๆ อาจอ่อน วิธีแก้คืออ่านซ้ำหรือขยายภาพ
\item ความรู้ภาษา (language prior คือสิ่งที่โมเดลคาดจากภาษาโดยไม่ดูภาพ) อาจทับหลักฐานในภาพ วิธีแก้คือหักความรู้ภาษาออกแบบ contrastive decoding \cite{leng2024vcd,favero2024m3id}
\item mark อาจไม่อยู่ใน output (ข้อความที่โมเดลถอดออกมา) เลย เพราะโมเดลข้ามบรรทัด หาข้อความไม่เจอ หรือติดวนซ้ำ (loop) วิธีแก้คือปรับการ decode หรือการหาข้อความ
\end{itemize}
output อย่างเดียวบอกไม่ได้ว่าข้างในโมเดลเกิดอะไรขึ้น เราจึงเรียกบรรทัดนั้นว่า ``ไม่ปรากฏใน output'' โดยไม่อ้างว่าโมเดลมองไม่เห็น

เราศึกษาเรื่องนี้บน ThaiOCRBench ด้วย Typhoon OCR 1.5 และ checkpoint แบบเปิดอีกเจ็ดตัว ในงานอ่านทั้งหน้าและอ่านข้อความสั้น เราพัฒนาการวิเคราะห์บนชุด calibration (ชุดสำหรับพัฒนาวิธี) แล้วยืนยันบนชุด locked ที่เปิดครั้งเดียว บทความนี้มีผลงานสี่ข้อ
\begin{enumerate}
\item คะแนน mark ที่ไม่ขึ้นกับลำดับการอ่าน (หัวข้อ~\ref{sec:metric}) ให้คะแนนสระบน สระล่าง และวรรณยุกต์รายบรรทัด ไม่ว่าบรรทัดอยู่ตรงไหนใน output ข้อความเกินถูกหักคะแนน สำเนาจากการวนซ้ำไม่ได้คะแนน
\item การวินิจฉัยว่า mark หายที่ใด (หัวข้อ~\ref{sec:where} และ~\ref{sec:locked}) บนชุด locked ทั้งแปดโมเดล บรรทัดที่ไม่ปรากฏใน output ทำให้เสีย mark มากกว่าการอ่าน mark ผิดหลังพยัญชนะที่อ่านถูก 4.7 ถึง 15.8 เท่า ถ้านับ mark ที่อ่านผิดทุกแบบ จะมากกว่า 2.0 ถึง 3.6 เท่า
\item การทดสอบ forced-choice ว่าโมเดลเลือก mark ใดเมื่อได้ข้อความที่ถูก (หัวข้อ~\ref{sec:where}) บนชุด calibration เมื่อได้ข้อความก่อนหน้าที่ถูก Typhoon และ Pathumma เลือกวรรณยุกต์ถูกราว 99\% โดยภาพช่วยเล็กน้อย ข้อนี้ไม่จริงทุกโมเดล Qwen3-VL-4B เลือก mark ถูกบ่อยเท่ากัน แต่ก็อ่านผิดบางตัวขณะเขียนหน้า
\item การประเมินการ decode ที่ไม่ต้องฝึกเพิ่ม (หัวข้อ~\ref{sec:decoding} และ~\ref{sec:locked}) บนชุด locked การตัด output ที่การวนซ้ำ (ทำแบบ offline) เพิ่มค่าประมาณ mark F1 ในทั้ง 13 เซลล์ที่ลงทะเบียน (เซลล์คือหนึ่งโมเดลบนหนึ่งงาน) และผ่านการทดสอบที่ลงทะเบียนใน 12 เซลล์ และเปลี่ยน output ที่จบเองเพียง 1 จาก 2,400 รายการ บน Typhoon (ชุด calibration) repetition penalty ลด F1 ของหน้าที่ไม่วนซ้ำ ส่วน contrast (การหักความรู้ภาษาตอนเลือก token) ที่เราเขียนใหม่ ลบวรรณยุกต์ที่ถูกออก
\end{enumerate}
ในโมเดลที่ศึกษา mark ส่วนใหญ่ที่โมเดล OCR ไทยทำหาย อยู่ในบรรทัดหรือช่วงข้อความที่ไม่ปรากฏใน output ไม่ใช่ mark ที่อ่านผิด แต่ยังไม่ทราบว่าข้อความนั้นไม่ปรากฏเพราะเหตุใด โมเดลอาจหาบรรทัดไม่พบ ไม่ได้เขียน หรือวนซ้ำ และยังตัดเหตุจากภาพ เช่น ความละเอียดที่ใช้จริงต่ำ ออกไม่ได้ เรายังไม่ได้ทดสอบว่าการเห็นวรรณยุกต์ชัดขึ้นจะช่วยหรือไม่

\section{งานวิจัยที่เกี่ยวข้อง}
\subsection{OCR ภาษาไทยด้วย VLM}
ThaiOCRBench \cite{nonesung2025thaiocrbench} ประเมิน VLM ทั้งปิดและเปิดบนงานเอกสารไทย รวมถึงอ่านทั้งหน้าและอ่านข้อความสั้น ตัวชี้วัดเป็นแบบ n-gram และ edit distance คือ BLEU, METEOR, F1 และ normalized Levenshtein similarity กับ ANLS สำหรับอ่านข้อความสั้น บทนำของ benchmark ระบุอักขระซ้อนว่าเป็นลักษณะของภาษาไทยที่ชุดข้อมูลแบบเน้นอังกฤษไม่ครอบคลุม แต่ไม่มีคะแนนใดแยก mark

ต่อมามี OCR ไทยแบบเปิดหลายตัว Typhoon OCR \cite{nonesung2026typhoonocr} รุ่น 1.5 เป็น full fine-tune ของ Qwen3-VL-2B \cite{bai2025qwen3vl} ผู้พัฒนาระบุว่า infographic ยังต้องปรับปรุง Pathumma-llm-vision 3.0 \cite{nectec2026pathumma} เป็น fine-tune ของ Qwen3-VL-2B อีกตัว wayu-paxa-ocr-zero \cite{pipatanakul2026synthetic} fine-tune PaddleOCR-VL \cite{cui2025paddleocrvl} ด้วยหน้าไทยสังเคราะห์ ทุกงานรายงานความแม่นยำรวม ส่วนเราถามว่า mark หายที่ใด

\subsection{เครื่องหมายกำกับใน OCR ภาษาอื่น}
อักษรที่มีเครื่องหมายกำกับ (diacritic) มากเป็นที่รู้กันว่ายากสำหรับ OCR มีงานแก้ข้อความหลัง OCR ภาษาเวียดนาม \cite{do2025reference} และงานประเมิน OCR อาหรับที่แยกข้อความที่มีเครื่องหมายกำกับเป็นอีกเงื่อนไข \cite{elsafty2025arabic} เราไม่พบงานที่ให้คะแนน mark ไทยแยกจากลำดับการอ่านและข้อความเกิน

\subsection{การวนซ้ำและการตกหล่น}
การ decode ที่เลือกความน่าจะเป็นสูงสุดทำให้ข้อความวนซ้ำ \cite{holtzman2020curious} โมเดล OCR เอกสารก็เป็นเช่นกัน งานที่มีอยู่จัดการต่างกัน
\begin{itemize}
\item Nougat รายงานหน้าที่พิมพ์ประโยคซ้ำ และหยุดสร้างข้อความเมื่อสัญญาณความแปรปรวนของ logit ลดลง \cite{blecher2024nougat}
\item olmOCR ลองใหม่ที่ temperature สูงขึ้น \cite{poznanski2025olmocr}
\item pipeline ของ PaddleOCR-VL ตัดข้อความซ้ำหลัง decode \cite{paddlex2025truncate}
\item โค้ด DeepSeek-OCR ที่เผยแพร่ห้าม n-gram ซ้ำในหน้าต่างหนึ่ง \cite{deepseekocr2025code,wei2025deepseekocr}
\item CC-OCR นับการวนซ้ำเป็นการหลอน (hallucination) แบบหนึ่งของ OCR \cite{yang2025ccocr}
\end{itemize}
งานเหล่านี้มุ่งตรวจจับหรือเลี่ยงการวนซ้ำ ส่วนเราวัดว่าแต่ละวิธีทำให้ mark ในหน้าที่ไม่วนซ้ำเสียเท่าใด หนึ่งในนั้นคือ repetition penalty ของ CTRL \cite{keskar2019ctrl} ที่ package OCR บางตัวเปิดเป็นค่าเริ่มต้น

\subsection{ความรู้ภาษาและ decoding ที่ไม่ต้องฝึกเพิ่ม}
decoding ที่ไม่ต้องฝึกเพิ่มลดการหลอนของ VLM โดยเชื่อสิ่งที่ภาพรองรับมากกว่าสิ่งที่ความรู้ภาษาอย่างเดียวทำนาย
\begin{itemize}
\item VCD \cite{leng2024vcd} เทียบการแจกแจงของ output จากภาพจริงกับภาพที่บิดเบือน
\item M3ID \cite{favero2024m3id} เพิ่มอิทธิพลของภาพเหนือความรู้ภาษา
\item PAI \cite{liu2024pai} เพิ่ม attention ต่อ token ภาพ และหัก logit ของอินพุตที่มีแต่ข้อความ
\end{itemize}
วิธีเหล่านี้สมมติว่าความรู้ภาษาเป็นปัญหา แต่ในข้อมูลของเรา ความรู้ภาษาเรื่องวรรณยุกต์ไทยมักถูก วิธีของเราเป็นการเขียนใหม่ จึงเรียกว่า M3ID-form และ PAI-style 

\section{คะแนน mark และการตั้งค่าการทดลอง}\label{sec:metric}\label{sec:setup}
mark ที่วัดคือ เครื่องหมายเหนือพยัญชนะหกตัว (\textit{กั กิ กี กึ กื ก็} คือสระห้าตัวกับไม้ไต่คู้) สระล่างสองตัว (\textit{กุ กู}) และวรรณยุกต์สี่ตัว (\textit{ก่ ก้ ก๊ ก๋}) หน้าเต็มมีเฉลยหลายบรรทัด และโมเดลอาจอ่านคอลัมน์ในลำดับต่างจากเฉลย การจัดแนวแบบ global จึงนับ mark ที่อ่านถูกเป็นข้อผิด

\subsection{นิยาม}
คะแนนมีสองขั้น
\begin{enumerate}
\item ขั้นแรกใช้บรรทัดเฉลยที่ยาวอย่างน้อยแปดตัวอักษร เรียงจากยาวไปสั้น แต่ละบรรทัดหาช่วงใน output ที่ตรงกันที่ character error rate ต่ำกว่า 40\% โดยข้ามอักขระที่ถูกจองแล้ว เราให้คะแนน mark ที่ตรงกัน แล้วจองช่วงนั้นไว้ ช่วงเดียวจึงให้คะแนนสองบรรทัดไม่ได้
\item ขั้นที่สอง บรรทัดที่เหลือจัดแนวแบบ global ตามลำดับเฉลยกับ output ที่ตัดอักขระที่จองออกแล้ว
\end{enumerate}
ขั้นที่สองจำเป็น เพราะหากไม่มี mark ในบรรทัดที่อ่านได้แย่จะหลุด และคะแนนไม่ผ่านการตรวจกับโมเดลฐาน (นิยามในย่อหน้าโมเดลด้านล่าง)

บรรทัดเฉลย $r$ ที่ยาวอย่างน้อยแปดตัวอักษร (Unicode code point) จับคู่กับช่วง $h$ ของ output เมื่อ
\begin{equation}
\mathrm{CER}(r,h)=\frac{S+D+I}{|r|}<0.40,\qquad |r|\ge 8.
\end{equation}
โดย $S$ $D$ และ $I$ คือจำนวนการแทนที่ การลบ และการเพิ่มของการจัดแนว การจัดหมวดในหัวข้อ~\ref{sec:where} ใช้ 0.20 แทน 0.40 เกณฑ์ทั้งสามค่า (40\% 20\% และแปดตัวอักษร) เรากำหนดบนชุด calibration และยังไม่ได้ทดสอบความไวต่อค่าเหล่านี้ บรรทัดที่อ่านแย่จนไม่ผ่านเกณฑ์ทั้งสองจึงถูกนับว่าหาย

mark ได้คะแนนเมื่อถูกจัดแนวกับ mark ตัวเดียวกันใน output ไม่ว่าพยัญชนะข้างๆ จะถูกหรือไม่ ให้ $G_i$ เป็นจำนวน mark ในเฉลยของข้อ $i$ ให้ $Y_i$ เป็นจำนวน mark ใน output เต็ม และให้ $C_i$ เป็นจำนวน mark ที่ได้คะแนน Recall คือ mark ที่ได้คะแนนหารด้วย mark ในเฉลย Precision คือ mark ที่ได้คะแนนหารด้วย mark ทั้งหมดใน output จึงลงโทษข้อความเกิน (ข้อความใน output ที่เฉลยไม่มี) F1 คำนวณจากจำนวนนับที่รวมทุกข้อ
\begin{equation}
R=\frac{\sum_i C_i}{\sum_i G_i},\quad
P=\frac{\sum_i C_i}{\sum_i Y_i},\quad
F_1=\frac{2\sum_i C_i}{\sum_i G_i+\sum_i Y_i}.
\end{equation}
ถ้าตารางหรือรูปไม่ระบุเป็นอื่น คะแนนใช้แบบมาตรฐานนี้ ส่วนแบบที่กันการวนซ้ำด้านล่างรายงานเฉพาะที่ระบุชื่อไว้

\subsection{แบบที่กันการวนซ้ำ}
แบบนี้คิดคะแนนหลังย่อแต่ละช่วงวนซ้ำให้เหลือสำเนาเดียว ช่วงวนซ้ำคือหน่วยยาว 4 ถึง 400 ตัวอักษรที่มีตัวอักษรอยู่ และซ้ำตรงตัวอย่างน้อยแปดครั้ง (หกครั้งหากยาว 50 ตัวอักษรขึ้นไป) ให้ $\rho(y)$ เป็น output $y$ ที่ย่อแต่ละช่วงวนซ้ำเหลือหนึ่งสำเนา คะแนนคิดจาก $\rho(y_i)$ แต่ตัวหารของ precision ยังนับ mark ทุกตัวใน output เต็ม สำเนาจึงถูกหักคะแนนและไม่ได้ recall
\begin{equation}
P_{\mathrm{RA}}=\frac{\sum_i C_i(\rho(y_i))}{\sum_i Y_i},\quad
R_{\mathrm{RA}}=\frac{\sum_i C_i(\rho(y_i))}{\sum_i G_i}.
\end{equation}

เราเพิ่มแบบนี้เพราะการจัดแนวจับสำเนา 434 ชุดของหน่วยซ้ำเดียวเข้ากับพยางค์ทั่วไปในเฉลย ถ้าไม่มีการวนซ้ำ สองคะแนนเท่ากัน

\subsection{การตรวจสอบ}
เราตรวจคะแนนด้วยสี่ข้อ
\begin{itemize}
\item การอ่านหน้าเดียวกันสองรอบต่อกันได้คะแนนต่ำกว่ารอบเดียวที่ดีกว่ามาก
\item recall ของโมเดลฐานต่างจากการจัดแนวแบบ global ไม่เกิน 0.3 จุดทุกเซลล์
\item recall หน้าเต็มของ Typhoon เพิ่ม 1.6 ถึง 1.8 จุดเมื่อไม่คิดลำดับ
\item การคิดลำดับอย่างเดียวลด mark F1 หน้าเต็มของ Typhoon จาก 94.9 เป็น 93.1 สำหรับ prompt BQ (95.3 เป็น 93.7 สำหรับ prompt TC ทั้งสอง prompt นิยามด้านล่าง) โมเดลฐานแทบไม่เปลี่ยน
\end{itemize}

\subsection{ข้อมูล}
เราใช้ชุดทดสอบ (test split) ของ ThaiOCRBench \cite{nonesung2025thaiocrbench} สองงาน คืออ่านทั้งหน้าและอ่านข้อความสั้น ชุด calibration สุ่มด้วย seed แบ่งชั้นตามงานและหมวด มีหน้าเต็ม 69 ข้อ และงานอ่านข้อความสั้น (text recognition) 109 ข้อ ส่วนที่เหลือเป็นชุด locked (ชุดที่ปิดไว้) มีหน้าเต็ม 128 ข้อ ข้อความสั้น 224 ข้อ

เราไม่เปิดชุด locked จนกฎทุกข้อในหัวข้อ~\ref{sec:where} และ~\ref{sec:decoding} ล็อกครบ และเปิดครั้งเดียวเพื่อยืนยันในหัวข้อ~\ref{sec:locked} เรารายงาน output และคะแนนรายข้อ โดยไม่เผยแพร่ภาพของ benchmark

\subsection{โมเดล}
เราประเมินแปด checkpoint
\begin{itemize}
\item Typhoon OCR 1.5 (2B) และโมเดลฐานของมัน คือ Qwen3-VL-2B (เรียกสั้นว่าโมเดลฐาน)
\item Pathumma-llm-vision 3.0
\item Qwen3-VL-4B และ Qwen-SEA-LION-v4-4B-VL \cite{aisingapore2026sealion}
\item DeepSeek-OCR-2 \cite{wei2026deepseekocr2}
\item โมเดลอ่านทีละกรอบขนาด 0.9B สองตัว คือ PaddleOCR-VL-1.6 และ wayu-paxa-ocr-zero (Wayu)
\end{itemize}
สองตัวสุดท้ายเรานำมาอ่านทั้งหน้าด้วย ซึ่งเกินที่ออกแบบไว้ หัวข้อ~\ref{sec:where} และ~\ref{sec:decoding} รายงานหกโมเดลที่ไม่รวม SEA-LION และ DeepSeek-OCR-2 สองตัวนี้เข้ามาในหัวข้อ~\ref{sec:locked}

DeepSeek-OCR-2 ใส่ GPU T4 แบบ fp32 ไม่ได้ ตัวเลขทั้งหมดของมันจึงเป็น fp16 จากโค้ดรันของโมเดลนี้เอง ทุก checkpoint ตรึง revision และเราใช้ revision ของ Typhoon ที่เป็น Apache-2.0 อย่างเดียว revision, seed และรหัสการรันบันทึกในคลังโค้ดของโครงการ

\subsection{Prompt และการ decode}
prompt ที่ใช้มีดังนี้
\begin{itemize}
\item BQ คือคำถามของ benchmark ในแต่ละข้อ
\item TC คือ prompt ที่ model card ของ Typhoon แนะนำ
\item PaddleOCR-VL และ Wayu ใช้ prompt \code{OCR:}
\item Pathumma และ DeepSeek-OCR-2 ใช้ prompt จาก model card ของตน
\end{itemize}
สำหรับ Typhoon BQ เป็น prompt หลัก TC เป็นรอง การ decode เป็นแบบ greedy ไม่เกิน 3,072 token ใหม่ ด้วย fp16 บน GPU Kaggle T4 แต่ละผลมาจากการรันครั้งเดียว output แบบ greedy กับที่ตัดที่การวนซ้ำมาจากการรันเดียวกัน จึงเป็นการเทียบแบบจับคู่ output ถูกแปลงเป็นข้อความธรรมดาก่อนให้คะแนน

\subsection{สถิติ}
ช่วงความเชื่อมั่นเป็น bootstrap แบบ percentile 95\% รายข้อ โดยข้อที่ใช้ภาพเดียวกันถูกสุ่มพร้อมกัน หัวข้อ~\ref{sec:where} และ~\ref{sec:decoding} มาจากชุด calibration จึงเป็นผลเบื้องต้น และกฎ decode บางข้อเขียนหลังเห็นชุดนี้ ซึ่งเราระบุในจุดนั้น หัวข้อ~\ref{sec:locked} ลงทะเบียนก่อนให้คะแนน output ชุด locked ใดๆ ใช้ bootstrap 10,000 รอบ และการแก้ Holm หนึ่งชุด การทดสอบ S1 และ R1 เป็น bootstrap ด้านเดียว S2 เป็นการทดสอบ binomial แบบแม่นยำด้านเดียวเทียบกับเกณฑ์ 1\% และค่า $p$ ทั้ง 22 ค่าถูกปรับด้วย Holm ที่ 0.05 S1 ผ่านเมื่อ $p$ ที่ปรับแล้วต่ำกว่า 0.05 และ R1 ผ่านเมื่อขอบล่างของช่วง 95\% ของอัตราส่วนยอดรวมสูงกว่า 3

\section{mark หายไปที่ใด}\label{sec:where}
\subsection{หกโมเดล}
บนชุด calibration Typhoon ได้ mark F1 94.9 บนหน้าเต็ม และ 81.7 บนข้อความสั้น (รูปที่~\ref{fig:models} และตารางที่~\ref{tab:models}) Pathumma-llm-vision 3.0 \cite{nectec2026pathumma} เป็นโมเดล fine-tune ภาษาไทยอีกตัวจากโมเดลฐานเดียวกัน และได้ใกล้เคียง คือ 89.8 และ 77.6 ด้วย prompt จาก model card

ส่วนต่างของ Pathumma จาก Typhoon คือ $-$5.1 จุดเปอร์เซ็นต์บนหน้าเต็ม (ช่วง $-$14.3 ถึง +1.6) และ $-$4.1 บนข้อความสั้น (ช่วง $-$15.1 ถึง +8.0) เราจึงบอกไม่ได้ว่าตัวใดดีกว่า Pathumma สูงกว่าโมเดลฐานมาก คือ 48.9 จุดบนหน้าเต็ม และ 53.3 จุดบนข้อความสั้น และ Typhoon สูงกว่านั้นอีก

อีกสามตัวต่ำกว่า Typhoon มาก ตัวดีที่สุดคือ Wayu ซึ่งต่ำกว่า 29.6 จุดบนหน้าเต็ม (ช่วง 19.2 ถึง 39.9) โมเดลฐานที่ใหญ่ขึ้นสองเท่าเปลี่ยน F1 หน้าเต็มเพียง +5.0 (ช่วง $-$4.0 ถึง +14.1) และอัตราผิดวรรณยุกต์ยังราว 9.4 ถึง 9.9\% Typhoon มี character error rate ต่อหน้าเป็นมัธยฐาน 2.8\%

การเทียบนี้มีข้อควรระวังสองข้อ
\begin{itemize}
\item โมเดลถูกเทียบตามที่เผยแพร่ ซึ่งมีจำนวน visual token ต่างกัน
\item สำหรับสองโมเดลที่ดีที่สุด เราตัดการซ้อนกับ benchmark ไม่ได้ model card ของ Pathumma รายงานคะแนน ThaiOCRBench โดยไม่ระบุว่าข้อมูลฝึก 377K ตัวอย่างคืออะไร และ Typhoon มาจากกลุ่มที่เผยแพร่ ThaiOCRBench
\end{itemize}
เราจึงเรียก Typhoon เพียงว่าได้คะแนนสูงที่สุดในที่เราวัด และใกล้เคียง Pathumma

ในงานอ่านข้อความสั้น สองโมเดลพลาดต่างกัน Pathumma พบ mark ที่ขอมากกว่า (recall สูงกว่า 14.3 จุด) แต่มักถอดทั้งภาพแทนช่วงที่ถาม precision จึงต่ำ

\begin{figure}[t]
\centering\includegraphics[width=\figwidth]{fig_models_\figlang.pdf}
\caption{Mark F1 บนชุด calibration (ก) อ่านทั้งหน้า (ข) อ่านข้อความสั้น ใช้ prompt BQ ยกเว้นโมเดล 0.9B สองตัวใช้ \code{OCR:} และ Pathumma ใช้ prompt จาก model card}\label{fig:models}
\end{figure}

\begin{table}[t]
\caption{Mark F1 อัตราผิดของวรรณยุกต์หลังอ่านพยัญชนะถูก และสัดส่วน output ที่ชนขีดจำกัด token (\%) Pathumma ใช้ prompt ของ model card}\label{tab:models}
\centering\footnotesize
\fit{\begin{tabular}{@{}l rrr rr@{}}
\toprule
 & \multicolumn{3}{c}{อ่านทั้งหน้า} & \multicolumn{2}{c}{อ่านข้อความสั้น}\\
\cmidrule(lr){2-4}\cmidrule(l){5-6}
โมเดล & F1 & ผิดวรรณยุกต์ & ชนขีดจำกัด & F1 & ชนขีดจำกัด\\
\midrule
Typhoon OCR 1.5 & \textbf{94.9} & 0.65 & 2.9 & \textbf{81.7} & 0.9\\
Pathumma 3.0 & 89.8 & 0.99 & 5.8 & 77.6 & 1.8\\
Wayu & 65.2 & 5.27 & 26.1 & 46.0 & 20.2\\
PaddleOCR-VL-1.6 & 53.3 & 14.57 & 39.1 & 47.9 & 14.7\\
Qwen3-VL-4B & 45.9 & 9.89 & 27.5 & 23.5 & 11.9\\
Qwen3-VL-2B & 41.0 & 9.41 & 43.5 & 24.3 & 14.7\\
\bottomrule
\end{tabular}}
\end{table}

\subsection{สำเนาของการวนซ้ำเคยได้ recall}
คะแนนแบบกันการวนซ้ำเปลี่ยน Typhoon ไม่เกิน 0.2 จุด (F1 94.9 เป็น 94.8, recall 96.1 เท่าเดิม) ส่วนโมเดลฐาน (BQ) หน้าเต็ม recall ลดจาก 65.0 เป็น 54.6 precision จาก 29.9 เป็น 25.1 และ F1 จาก 41.0 เป็น 34.4 ราวหนึ่งในหกของ mark ที่คะแนนแบบไม่ขึ้นกับลำดับให้โมเดลฐานมาจากสำเนาวนซ้ำ

เราอ่านทุกช่วงที่ถูกตัด ทั้ง 126 ช่วงชนขีดจำกัด token ช่วงสั้นสุดมี 15 สำเนา และไม่มีข้อความปกติถูกตัด

\subsection{mark ที่หายส่วนใหญ่ไม่ปรากฏใน output}
บนหน้าเต็ม Typhoon ทำ mark ผิด 6.2\% (BQ) และ 7.1\% (TC) ของ mark ในเฉลย เราจัดข้อผิดแต่ละตัวเข้าหนึ่งในสี่หมวดโดยอัตโนมัติ (รูปที่~\ref{fig:causes}) ตารางที่~\ref{tab:causes} แสดงสัดส่วนของ mark ในเฉลยที่แต่ละหมวดคิดเป็น

\begin{table}[t]
\caption{Typhoon หน้าเต็ม สัดส่วนของ mark ในเฉลย (\%) แยกตามหมวด สำหรับ prompt BQ และ TC}\label{tab:causes}
\centering\footnotesize
\begin{tabular}{@{}l rr@{}}
\toprule
หมวด & BQ & TC\\
\midrule
ทั้งบรรทัดหาย & 2.3 & 2.9\\
ข้อความขาดเป็นช่วงในบรรทัดที่ยังอยู่ & 0.7 & 0.8\\
mark ถูกอ่านผิด & 1.5 & 2.4\\
\quad ผิดขณะพยัญชนะอ่านถูก & 0.4 & 0.6\\
\quad ผิดพร้อมพยัญชนะ & 1.1 & 1.8\\
\midrule
ข้อผิด ไม่นับลำดับการอ่าน & 4.6 & 6.1\\
อ่านในลำดับอื่น (ไม่ใช่ข้อผิด) & 1.6 & 1.0\\
\bottomrule
\end{tabular}
\end{table}

``อ่านในลำดับอื่น'' หมายถึงมีช่วงใน output ตรงกับบรรทัดที่ character error rate ต่ำกว่า 20\% (สำหรับบรรทัดยาวอย่างน้อยแปดตัวอักษร) ส่วนนี้คือความต่างของลำดับการอ่าน ไม่ใช่ข้อผิดของการอ่าน ข้อผิดส่วนใหญ่เป็นบรรทัดที่หายหรือขาดเป็นช่วง

ใน output จริง บรรทัดที่หายส่วนใหญ่เป็นข้อความใน infographic และกราฟ โฆษณาที่แทรก และชื่อเรื่องหรือส่วนประกอบหน้า ซึ่ง Typhoon เขียนตัวแทนภาพไว้แทน สอดคล้องกับที่ผู้พัฒนาระบุเรื่อง infographic \cite{nonesung2026typhoonocr}

การจัดหมวดเป็นแบบอัตโนมัติ และเราดูภาพเฉพาะหกหน้าที่มี mark ที่หายเป็นส่วนใหญ่ บรรทัดที่อ่านแย่จนเกินเกณฑ์ character error rate จึงถูกนับว่าหาย

รูปแบบเดียวกันพบในอีกห้าโมเดล (รูปที่~\ref{fig:causes}) ที่นี่เราจัดหมวดบน output ที่ตัดที่การวนซ้ำครั้งแรก เพราะ output ที่วนซ้ำยาวกว่าเฉลยหลายเท่า จึงเชื่อการจัดแนวไม่ได้
\begin{itemize}
\item บรรทัดที่หายรวมช่วงที่ขาดทำให้เสีย mark 3.2\% ใน Typhoon, 9.3\% ใน Pathumma และ 29 ถึง 44\% ในอีกสี่โมเดล
\item mark ผิดหลังพยัญชนะที่อ่านถูกเสีย 0.4\% ใน Typhoon, 0.8\% ใน Pathumma และ 2.7 ถึง 4.3\% ในอีกสี่โมเดล
\end{itemize}
รูปที่~\ref{fig:causes} รวมการอ่านผิดทั้งหมดเป็นส่วนเดียว ซึ่งรวม mark ที่ผิดพร้อมพยัญชนะด้วย แต่การเทียบต่อไปนี้และ R1 ในหัวข้อ~\ref{sec:locked} ใช้เฉพาะ mark ผิดหลังพยัญชนะที่อ่านถูก เมื่อคำนวณทีละโมเดล บรรทัดที่หายจึงทำให้เสีย mark มากกว่าการอ่านผิด 7 ถึง 12 เท่าในทุกโมเดล

สิ่งที่วัดคือการไม่ปรากฏใน output ซึ่งไม่เท่ากับโมเดลอ่านไม่ได้ เพราะโมเดลอาจเห็นบรรทัดแต่ไม่เขียนออกมา การแยกสองกรณีต้องมีการทดสอบ เช่น ให้อ่านบรรทัดนั้นเดี่ยวๆ จากภาพที่ตัดมา ซึ่งเรายังไม่ได้ทำกับหน้าเต็ม ในโมเดลที่วนซ้ำ บรรทัดหลังการวนซ้ำนับว่าหาย และเราแยกไม่ได้ว่าหายเพราะการวนซ้ำหรือเพราะโมเดลไม่อ่านอยู่แล้ว

\begin{figure}[t]
\centering\includegraphics[width=\figwidth]{fig_causes_\figlang.pdf}
\caption{mark หน้าเต็มที่หาย แยกตามหมวด เป็นสัดส่วนของ mark ในเฉลยทั้งหมด สำหรับ output ที่ตัดที่การวนซ้ำครั้งแรก ส่วนอ่านผิดรวม mark ที่ผิดพร้อมพยัญชนะ ยอดรวมแต่ละแท่งไม่รวมส่วนสีเทา ซึ่งเป็นความต่างของลำดับการอ่านและไม่ใช่ข้อผิด}\label{fig:causes}
\end{figure}

\subsection{เมื่อได้บริบทที่ถูก วรรณยุกต์แทบถูกเสมอ}
เราทดสอบแบบ forced-choice คือให้โมเดลเลือกจากตัวเลือกที่กำหนด เราป้อนข้อความที่ถูกต้องจนถึงตำแหน่งของ mark แล้วให้โมเดลให้คะแนนรูปแบบ mark ที่เป็นไปได้ ณ ตำแหน่งนั้น ทั้งมีและไม่มีภาพ (คำนวณด้วย fp32 โดย tokenize ในบริบท) บนหน้าเต็มของชุด calibration Typhoon เลือกวรรณยุกต์ถูก 99.7\% เมื่อมีภาพ และ 98.8\% เมื่อไม่มี โมเดลฐานได้ 98.1\% และ 98.9\% ภาพจึงช่วย Typhoon 0.9 จุด แต่ลดผลของโมเดลฐานลง 0.8 จุด (รูปที่~\ref{fig:prior}ก)

การทดสอบนี้ให้บริบทที่ถูก จึงต่างจากการให้โมเดลเขียนเอง ผลแสดงสิ่งที่โมเดลทำนายจากบริบทนั้น และเพราะภาพเพิ่มแค่ราวหนึ่งจุด จึงไม่ได้แสดงว่าช่องทางภาพแยกวรรณยุกต์ได้ ตัวเลือกคือรูปแบบที่เป็นไปได้ของ mark โดยไม่มี negative ที่ยากเป็นพิเศษ

\subsection{อีกสองโมเดล}
ผลของการทดสอบ forced-choice
\begin{itemize}
\item Pathumma ได้ 99.7\% เมื่อมีภาพ 98.7\% เมื่อไม่มี (เพิ่ม +1.0) และ 94.8\% เมื่อหักคะแนนแบบไม่มีภาพที่ $\lambda=0.5$
\item Qwen3-VL-4B เลือก mark ถูกบ่อยเท่ากัน (98.7\% เมื่อมีภาพ) และภาพไม่ช่วย (99.1\% เมื่อไม่มี)
\end{itemize}
ทั้งสองต่างกันเมื่อเขียนหน้าเอง ณ ตำแหน่งเดียวกับการทดสอบ
\begin{itemize}
\item ที่ตำแหน่งซึ่ง Pathumma อ่านพยัญชนะฐานถูก วรรณยุกต์ถูก 99.4\% เทียบกับ 99.8\% ในการทดสอบ
\item Qwen3-VL-4B ได้ 91.3\% เทียบกับ 99.1\% คืออ่านวรรณยุกต์ผิดราวหนึ่งในสิบเอ็ดตัวขณะเขียนเอง ทั้งที่การทดสอบแสดงว่าเลือกตัวที่ถูกได้
\end{itemize}
สองตัวเลขนี้ต่างกันที่ความแม่นยำเชิงตัวเลขด้วย (เขียนเองด้วย fp16 ทดสอบด้วย fp32) ส่วนต่างจึงปนผลของการ decode กับความแม่นยำ ภาพช่วยเล็กน้อยในโมเดล fine-tune ภาษาไทยสองตัว และไม่ช่วยโมเดลทั่วไป เราไม่ขยายผลเกินสี่โมเดลนี้

สำหรับ Qwen3-VL-4B ข้อที่ว่า mark ส่วนใหญ่หายจาก output มากกว่าถูกอ่านผิดยังจริงในทิศทาง (บน calibration ข้อความที่หายทำให้เสีย mark มากกว่าการอ่านผิดเหล่านี้ 7.3 เท่า) แต่ไม่จริงในความหมายว่า mark แทบไม่เคยถูกอ่านผิด

\begin{figure}[t]
\centering\includegraphics[width=\figwidth]{fig_prior_\figlang.pdf}
\caption{(ก) ความแม่นยำวรรณยุกต์ที่เพิ่มเมื่อมีภาพ ณ ตำแหน่งที่ลงทะเบียน (ชุด calibration) (ข) ความแม่นยำวรรณยุกต์ของ Typhoon เมื่อหักคะแนนแบบไม่มีภาพด้วยน้ำหนัก $\lambda$}\label{fig:prior}
\end{figure}

\section{วิธี decode ที่ไม่ต้องฝึกโมเดลเพิ่ม และผลต่อ mark}\label{sec:decoding}
\subsection{การวนซ้ำพบน้อยแต่เสียหายมาก}
output ของ Typhoon 356 รายการ (สอง prompt สองงาน) วนซ้ำ 11 รายการ ในแต่ละเซลล์ หน้าที่วนซ้ำสองถึงห้าหน้ามี mark เกินรวม 68 ถึง 76\% ของทั้งหมด หากไม่มีหน้าเหล่านี้ F1 หน้าเต็ม (BQ) จะเป็น 97.4 แทน 94.9 ตัวเลขนี้คือ F1 ของหน้าที่เหลือ ไม่ได้เป็นผลของวิธีใด

\subsection{สามวิธีจัดการการวนซ้ำ}
เราเทียบสามวิธีกับ greedy decoding บน Typhoon (รูปที่~\ref{fig:decoders} หน้าเต็ม BQ)
\begin{itemize}
\item Repetition penalty 1.1 ค่าที่ package ของ Typhoon ใช้ เป็น penalty ของ CTRL \cite{keskar2019ctrl}
\item n-gram blocking แบบโค้ด DeepSeek-OCR ห้าม 30-gram ซ้ำภายใน 90 token
\item หยุดเมื่อพบการวนซ้ำ ถือว่าพบเมื่อหน่วยหนึ่งซ้ำแปดครั้ง (หกครั้งหากยาว 50 ตัวอักษรขึ้นไป) จากนั้นหยุด decode และตัด output ย้อนกลับไปที่ท้ายสำเนาแรก greedy decoding ให้ผลเดิมเสมอเมื่อ prefix เดิม เราจึงคำนวณจุดตัดเดียวกันแบบ offline จาก output greedy ที่ครบแล้ว และตรวจตัวที่หยุดตอน decode จริงเทียบกับจุดตัดนี้ในหัวข้อ~\ref{sec:locked}
\end{itemize}

\begin{figure}[t]
\centering\includegraphics[width=\figwidth]{fig_decoders_\figlang.pdf}
\caption{การเปลี่ยนของ mark F1 หน้าเต็มเทียบกับ greedy decoding (Typhoon, BQ) แท่งสีน้ำเงินคือทุกหน้า พร้อมช่วง 95\% แท่งสีส้มคือหน้าที่ไม่วนซ้ำ}\label{fig:decoders}
\end{figure}

\subsection{Repetition penalty}
เปลี่ยน output เกือบทุกรายการ ไม่มี output หน้าเต็ม BQ ที่ไม่วนซ้ำรายการใดเหมือน greedy F1 ของหน้าที่ไม่วนซ้ำลด 1.1 ถึง 1.3 จุด ในข้อความสั้น (BQ) อัตราผิดวรรณยุกต์เพิ่มจาก 17.1\% เป็น 22.6\% เมื่อรวมทุกหน้า ผลแยกจากศูนย์ไม่ได้ (+0.89 ช่วง $-$2.70 ถึง +4.81)

ข้อความไทยมีพยางค์และวรรณยุกต์ซ้ำได้ตามปกติ penalty ทั้งหน้าอาจกดสิ่งเหล่านี้ เรายังไม่ได้ทดสอบคำอธิบายนี้

\subsection{n-gram blocking}
เมื่อไม่มีการวนซ้ำ output 92 ถึง 100\% ไม่เปลี่ยน และ F1 เปลี่ยนไม่เกิน 0.05 วิธีนี้หยุดการวนซ้ำได้ 6 จาก 11 รายการ

ในรายการที่เหลือ การวนซ้ำไม่หยุด โมเดลเปลี่ยนไปวนซ้ำข้อความที่เกือบเหมือนเดิมแทน เช่น ชื่อซอยที่เพี้ยนเป็นชื่อคล้ายกันสามชื่อ ข้อความแบบนี้มี mark มากกว่าเดิม จึงมี mark ส่วนเกินเพิ่ม และ F1 หน้าเต็ม (BQ) ต่ำลง 0.68 เราตีความว่าการวนซ้ำเป็น ``หลุมดึงดูด'' ของ decoder (attractor) เมื่อห้ามซ้ำแบบตรงตัว โมเดลจึงไปติดหลุมข้างเคียง นี่เป็นการตีความที่ยังไม่ได้ทดสอบ

\subsection{หยุดเมื่อพบการวนซ้ำ}
ไม่เปลี่ยน output ที่ไม่วนซ้ำเลยในทุกเซลล์ mark F1 หน้าเต็มเพิ่ม 0.84 (BQ) และ 0.67 (TC) และ F1 ข้อความสั้นเพิ่ม 3.30 และ 2.39 การหยุดยังประหยัด token ที่สร้าง 11.2\% recall ลดไม่เกิน 0.16 จุด เพราะการหยุดกู้ข้อความที่การวนซ้ำข้ามไปไม่ได้

เราเขียนกฎนี้หลังเห็นการเทียบก่อนหน้า แต่ละเซลล์มี output วนซ้ำเพียงไม่กี่รายการ ทุกช่วงจึงมีศูนย์อยู่ในช่วง

\subsection{การหยุดกับโมเดลอื่น}
โมเดลอื่นวนซ้ำบ่อยกว่า Typhoon มาก บนหน้าเต็ม โมเดลฐาน Qwen3-VL-4B PaddleOCR-VL และ Wayu วนซ้ำ 18 ถึง 30 จาก 69 รายการ เทียบกับ 2 ของ Typhoon และ 4 ของ Pathumma

การหยุดแบบ offline เดียวกันบน output greedy ของโมเดลเหล่านี้เพิ่ม mark F1 ทุกโมเดลทั้งสองงาน ไม่มีช่วงใดต่ำกว่าศูนย์ และเจ็ดจากสิบสองเซลล์ไม่มีศูนย์อยู่ในช่วง บนหน้าเต็ม ผลเพิ่มมากเมื่อโมเดลวนซ้ำบ่อย เช่น +24.4 ในโมเดลฐาน และ +10.8 ใน Wayu (ช่วง 5.3 ถึง 16.1) และน้อยเมื่อไม่ค่อยวนซ้ำ คือ +0.8 ใน Typhoon และ +3.0 ใน Pathumma หัวข้อ~\ref{sec:locked} ทดสอบซ้ำบนข้อที่ไม่เคยเห็นกับทั้งแปดโมเดล

ผลเพิ่มมาจาก precision โดยแลกกับ recall ที่ลดสูงสุด 10.6 จุดบนหน้าเต็ม ซึ่งเป็นไปตามที่คะแนนบ่งไว้ เพราะ precision นับ mark ทุกตัวใน output recall ที่ลดส่วนหนึ่งเป็นคะแนนที่สำเนาวนซ้ำเคยได้ ที่เหลือคือข้อความหลังการวนซ้ำที่ถูกตัด

ในสิบสองเซลล์ การหยุดเปลี่ยน output ที่จบเองเพียง 3 จาก 898 รายการ สองรายการเป็นของ PaddleOCR-VL ที่ออกจากการวนซ้ำแบบเสื่อมได้เอง แล้วการตัดทำให้ข้อความจริงที่ตามมาหายไป repetition penalty 1.05 อย่างเดียว ซึ่งเป็นค่าในตัวอย่างของ Wayu เอง ไม่ช่วย (+2.8 ช่วง $-$3.6 ถึง +9.6 บนหน้าเต็ม และ $-$11.3 ช่วง $-$22.5 ถึง +4.5 บนข้อความสั้น)

การหยุดประหยัด token ที่สร้าง 34\% บนหน้าเต็มของ Wayu ในการรันนี้ 72\% ของช่องว่าง recall หน้าเต็มระหว่าง Wayu กับ Typhoon อยู่ใน 18 หน้าที่ Wayu วนซ้ำ และไม่มีวิธี decode ใดที่เราทดสอบกู้คืนได้ เราบอกไม่ได้ว่าการวนซ้ำเป็นสาเหตุหรืออาการ

\subsection{การหักความรู้ภาษาออก}
ในการทดสอบ forced-choice การหักคะแนนแบบไม่มีภาพด้วยน้ำหนัก $\lambda$ ลดความแม่นยำวรรณยุกต์ของ Typhoon จาก 99.7\% เป็น 97.3\% ที่ $\lambda=0.5$ และ 42.3\% ที่ $\lambda=1$ (รูปที่~\ref{fig:prior}ข) โมเดลฐานลดเหลือ 84.7\% ที่ 0.5

ในการสร้างข้อความเองบนกลุ่มย่อย 24 ข้อ contrast แบบ M3ID-form ของเรา (หักคะแนนแบบไม่มีภาพออกจากคะแนนแบบมีภาพ) หยุดการวนซ้ำของโมเดลฐานได้ (character error rate 1.12 เป็น 0.21) แต่เพิ่มอัตราผิดวรรณยุกต์ในข้อที่เก็บไว้จาก 3.7\% เป็น 7.6\% ส่วนใหญ่ด้วยการลบวรรณยุกต์ (21 เป็น 55 ใน 12 ข้อ) บน Typhoon การเปลี่ยนแปลงเล็กและกระจุกในไม่กี่ข้อ

ผลทั้งสองสอดคล้องกับความรู้ภาษาเรื่องวรรณยุกต์ที่ถูกอยู่แล้วใน Typhoon และโมเดลฐาน ผลจากการสร้างข้อความเองอาศัยเพียง 24 ข้อและวิธีที่เขียนใหม่

\section{การยืนยันผลบนชุด locked}\label{sec:locked}
ทุกอย่างข้างต้นพัฒนาบนชุด calibration จากนั้นเราลงทะเบียนการทดสอบสามข้อพร้อมวิธีอ่านผลล่วงหน้า แล้วเปิดชุด locked ครั้งเดียว คือ 352 ข้อ (หน้าเต็ม 128 ข้อ ข้อความสั้น 224 ข้อ) กับแปดโมเดล คือหกโมเดลข้างต้นกับ Qwen-SEA-LION-v4-4B-VL และ DeepSeek-OCR-2 prompt, greedy decoding และ fp16 เป็นตามหัวข้อ~\ref{sec:setup} และการตัดที่การวนซ้ำครั้งแรกทำแบบ offline บน output greedy ของการรันเดียวกัน การทดสอบมีดังนี้
\begin{itemize}
\item S1 การหยุดช่วยไหม: ในแต่ละคู่โมเดล-งานที่วนซ้ำอย่างน้อย 5\% ของข้อ calibration การหยุดต้องเพิ่ม mark F1 มี 13 เซลล์เช่นนี้ บน calibration สัดส่วนข้อที่วนซ้ำเท่ากับสัดส่วนที่ชนเพดาน token ในตารางที่~\ref{tab:models} ทุกโมเดลและทุกงาน
\item S2 การหยุดทำลายของปกติไหม: การหยุดต้องเปลี่ยน output ที่จบเองไม่เกิน 1\% เมื่อรวมทุกโมเดลและทุกงาน
\item R1 mark ส่วนใหญ่หายเพราะข้อความไม่ปรากฏไหม: ในแต่ละโมเดลบนหน้าเต็ม mark ที่เสียกับบรรทัดที่หายและช่วงที่ขาดต้องเกินสามเท่าของที่เสียกับการอ่านผิดหลังพยัญชนะที่อ่านถูก โดยขอบล่างของช่วงความเชื่อมั่น 95\% สูงกว่าสาม
\end{itemize}
สำหรับโมเดล $m$ ให้ $L_{mi}$ และ $M_{mi}$ เป็น mark ในเฉลยของข้อ $i$ ที่อยู่ในบรรทัดที่หายและในช่วงที่ขาด และ $E_{mi}$ เป็น mark ที่ผิดหลังพยัญชนะที่อ่านถูก นับทั้งหมดบน output ที่ตัดที่การวนซ้ำครั้งแรก R1 คืออัตราส่วนของยอดรวม
\begin{equation}
R_m=\frac{\sum_i (L_{mi}+M_{mi})}{\sum_i E_{mi}}.
\end{equation}
การแก้ Holm หนึ่งชุด (การปรับค่า $p$ เมื่อทดสอบหลายข้อพร้อมกัน) ครอบคลุมทั้ง 22 การทดสอบ (S1 13 ข้อ S2 หนึ่งข้อ R1 แปดข้อ) ช่วงเป็น bootstrap รายหน้า 10,000 รอบ

\subsection{การหยุด}
mark F1 เพิ่มในทั้ง 13 เซลล์ที่เข้าเกณฑ์ และ 12 เซลล์ผ่านการทดสอบที่ลงทะเบียน (รูปที่~\ref{fig:locked} และตารางที่~\ref{tab:locked}) ใน 12 เซลล์นั้นเพิ่ม 4.8 ถึง 35.6 จุด โดย $p$ ที่ปรับแล้วไม่เกิน 0.0036

เซลล์ที่สิบสามคือ Pathumma หน้าเต็ม ไม่ผ่าน ผลเพิ่มเพียง +1.3 จุด (ช่วง 0.0 ถึง +3.5, $p$ ที่ปรับแล้ว 0.13) เซลล์นี้เข้าเกณฑ์บน calibration แบบเฉียดเกณฑ์ (วนซ้ำ 4 จาก 69 หน้า) บนชุด locked วนซ้ำ 4 จาก 128 หน้า (3.1\%) ซึ่งต่ำกว่าเกณฑ์ของการเข้าทดสอบ เมื่อโมเดลแทบไม่วนซ้ำ ผลเพิ่มย่อมเล็ก และขนาดตัวอย่างนี้ยืนยันไม่ได้ เกณฑ์ที่ลงทะเบียนต้องการให้ทุกเซลล์ที่เข้าเกณฑ์ผ่าน S1 จึงไม่ผ่านตามที่ลงทะเบียน และเราระบุเซลล์นี้ไว้

S2 ผ่าน การหยุดเปลี่ยน output ที่จบเอง 1 จาก 2,400 รายการ (0.04\% ขอบบน 95\% คือ 0.20\%) คือ output ของ Typhoon ที่ออกจากการวนซ้ำจริงได้เอง

\subsection{mark หายที่ไหน}
R1 ผ่านทั้งแปดโมเดล (ตารางที่~\ref{tab:locked}) บนหน้าเต็ม ข้อความที่หายทำให้เสีย mark มากกว่าการอ่านผิดหลังพยัญชนะที่อ่านถูก 4.7 เท่า (PaddleOCR-VL ช่วง 3.2 ถึง 6.7) ถึง 15.8 เท่า (Pathumma) ขอบล่างของ PaddleOCR-VL ใกล้เกณฑ์ที่สุด

ข้อนี้กล่าวถึงแปดโมเดลนี้บน benchmark เดียว

\begin{figure}[t]
\centering\includegraphics[width=\figwidth]{fig_locked_\figlang.pdf}
\caption{Mark F1 บนชุด locked เมื่อไม่หยุดและหยุดที่การวนซ้ำ (ก) อ่านทั้งหน้า (ข) อ่านข้อความสั้น prompt ตามที่ลงทะเบียน DeepSeek-OCR-2 เป็น fp16 เท่านั้น}\label{fig:locked}
\end{figure}

\begin{table*}[t]
\caption{ชุด locked การเปลี่ยนของ mark F1 เมื่อหยุด (จุด ช่วง 95\%) และ R1 คืออัตราส่วนของ mark ที่เสียกับข้อความที่หาย ต่อ mark ที่เสียกับการอ่านผิดหลังพยัญชนะที่อ่านถูก (หน้าเต็ม พร้อมช่วง 95\%) คอลัมน์สุดท้ายคืออัตราส่วนเดียวกันเมื่อนับ mark ที่อ่านผิดทุกแบบเป็นตัวหาร (ค่าจุดเดียว ไม่อยู่ในการทดสอบที่ลงทะเบียน) \textsuperscript{$\dagger$}ไม่ได้ทดสอบ เพราะโมเดลวนซ้ำน้อยกว่า 5\% ของข้อ calibration \textsuperscript{$\ddagger$}ไม่ผ่านการทดสอบที่ลงทะเบียน}\label{tab:locked}
\centering\footnotesize
\begin{tabular}{@{}l llll@{}}
\toprule
โมเดล & อ่านทั้งหน้า $\Delta$F1 & อ่านข้อความสั้น $\Delta$F1 & R1 & R1 นับอ่านผิดทุกแบบ\\
\midrule
Typhoon OCR 1.5 & +1.5 [$-$0.5, +4.1]\textsuperscript{$\dagger$} & +3.5 [0.0, +11.0]\textsuperscript{$\dagger$} & 8.2 [5.5, 14.3] & 2.7\\
Pathumma 3.0 & +1.3 [0.0, +3.5]\textsuperscript{$\ddagger$} & +3.0 [0.0, +6.7]\textsuperscript{$\dagger$} & 15.8 [7.5, 32.0] & 3.3\\
Wayu & +4.8 [+1.8, +8.3] & +29.1 [+18.7, +38.0] & 10.8 [7.1, 16.0] & 3.0\\
PaddleOCR-VL-1.6 & +8.4 [+4.5, +12.0] & +8.8 [+2.2, +14.6] & 4.7 [3.2, 6.7] & 2.0\\
Qwen3-VL-4B & +16.3 [+11.7, +20.3] & +7.0 [+4.2, +10.1] & 6.1 [3.7, 9.2] & 2.2\\
Qwen3-VL-2B & +18.4 [+14.3, +22.2] & +35.6 [+29.6, +40.9] & 7.8 [5.0, 11.7] & 3.6\\
SEA-LION v4 4B & +15.2 [+11.9, +18.4] & +6.9 [+3.8, +10.0] & 7.8 [5.6, 10.4] & 2.3\\
DeepSeek-OCR-2 (fp16) & +17.1 [+12.8, +21.5] & +30.0 [+21.8, +36.3] & 7.8 [5.8, 10.5] & 2.1\\
\bottomrule
\end{tabular}
\end{table*}

\subsection{การจัดอันดับ}
mark F1 ของ Typhoon สูงกว่าหกโมเดลที่ไม่ใช่ Pathumma ทั้งสองงาน โดยช่วงอยู่เหนือศูนย์ เทียบกับ Pathumma ไม่พบความต่าง

ด้วย prompt จาก model card ของ Pathumma ส่วนต่างคือ +1.5 จุดบนหน้าเต็ม (ช่วง $-$3.1 ถึง +6.5) และ +9.0 บนข้อความสั้น (ช่วง $-$3.6 ถึง +21.3) ที่ 128 และ 224 ข้อ ครึ่งความกว้างของช่วงราว 5 และ 12 จุด จึงไม่ใช่ความเท่ากัน และเราไม่เรียกตัวใดว่าดีที่สุด

การให้คะแนนแกะคำตอบในรูป JSON ของ Typhoon และสำหรับ output ที่จบเอง เก็บข้อความหน้าหลังแท็ก \code{<figure>} ที่ไม่ถูกปิด สองกฎนี้เพิ่มหลังให้คะแนน output ชุด locked ครั้งแรกแล้ว จึงเป็นการทำภายหลัง (post hoc) กฎนี้ใช้กับทุกโมเดล บนชุด locked มีเพียง Pathumma ที่มี output แบบนี้ และตัวเลขข้างต้นเท่าเดิมไม่ว่าใช้กฎนี้หรือไม่

F1 ชุด locked ของ Typhoon (92.6 และ 73.7) ต่ำกว่า F1 บน calibration (94.9 และ 81.7) เราไม่ได้ปรับอะไรกับข้อ locked ค่าที่ลดลงสอดคล้องกับการที่กฎเขียนจากชุด calibration แต่เราไม่ได้แยกสาเหตุ และใช้ตัวเลขชุด locked เป็นตัวเลขหลัก

\subsection{การหยุดตอนสร้างข้อความจริง}
output ชุด locked ถูกตัดแบบ offline เพื่อตรวจกฎตอน decode เราสร้างข้อความบนข้อ calibration ใหม่โดยต่อตัวหยุดเข้ากับการสร้างข้อความของทั้งแปดโมเดล คือ 185 ข้อที่การตัดแบบ offline เปลี่ยน output และ 72 ข้อที่สุ่มจากข้อที่การตัดไม่แตะ

output ที่หยุดตรงกับการตัด offline ทุกตัวอักษรใน 185 จาก 185 ข้อ และกฎไม่ทำงานใน 72 จาก 72 output ที่การตัดไม่แตะ output ที่วนซ้ำใช้ token 15.3\% ของเดิม การตรวจนี้ครอบคลุมเฉพาะข้อ calibration และเรายังไม่ได้วัดเวลาที่ประหยัด

\section{อภิปรายผลและข้อจำกัด}
เราเทียบการหยุดกับ repetition penalty และ n-gram blocking เฉพาะบน Typhoon ชุด calibration จึงไม่อ้างว่าดีกว่าวิธีอื่น ภาพนี้ไม่เหมือนกันทุกโมเดล Qwen3-VL-4B ยังอ่าน mark ผิดขณะเขียนหน้าด้วย

\subsection{ขอบเขต}
ข้อจำกัดด้านขอบเขตมีดังนี้
\begin{itemize}
\item งานนี้ใช้ benchmark เดียวและสองงาน
\item โมเดลเป็น checkpoint แบบเปิดตามที่เผยแพร่ ซึ่งมีจำนวน visual token และความละเอียดอินพุตต่างกันโดยเราไม่ได้ควบคุม
\item ห้าในแปดตัวไม่ได้ฝึกมาเฉพาะสำหรับ OCR ไทย และหลายตัวได้คะแนนต่ำ ข้อความที่ไม่ปรากฏมากกว่าการอ่านผิดจึงไม่น่าแปลกใจสำหรับตัวเหล่านั้น
\item โมเดลอ่านทีละกรอบสองตัวถูกใช้กับหน้าเต็มเกินที่ออกแบบ และ DeepSeek-OCR-2 รัน fp16 อย่างเดียวด้วยโค้ดรันของตัวเอง
\item Typhoon มาจากกลุ่มที่เผยแพร่ benchmark และข้อมูลฝึกของ Pathumma ไม่ถูกระบุ เราจึงตัดการซ้อนของข้อมูลไม่ได้และไม่จัดอันดับสองตัวนี้
\item Typhoon และ Pathumma วนซ้ำน้อย ช่วงส่วนใหญ่จึงมีศูนย์อยู่ในช่วง เราจึงถือผลของสองโมเดลนี้เป็นรูปแบบที่เห็นใน output ไม่ใช่ค่าประมาณขนาดผล
\item เรายังไม่มี control (กลุ่มเทียบ) สำหรับอักขระที่ไม่ใช่ mark ข้อความที่หายทำให้ทุกอักขระของบรรทัดหายไป เราจึงไม่ได้แสดงว่า mark หายในอัตราที่สูงกว่าพยัญชนะ
\end{itemize}

\subsection{การออกแบบ}
หัวข้อ~\ref{sec:where} และ~\ref{sec:decoding} มาจากชุด calibration และกฎการหยุด คะแนนแบบกันการวนซ้ำ และการให้คะแนน เขียนหลังเห็นข้อมูลชุดนั้น

หัวข้อ~\ref{sec:locked} คือการยืนยัน ลงทะเบียนก่อนให้คะแนน output ชุด locked ใดๆ และเปิดชุดนี้ครั้งเดียว ตัวเลขที่ควรอ้างมาจากหัวข้อนี้ ผลสนับสนุนการวินิจฉัยและการหยุดสำหรับแปดโมเดลนี้ และเราไม่อ้างเกินนั้น

ข้อจำกัดของการออกแบบ S1 ตามที่ลงทะเบียนไม่ผ่านเพราะหนึ่งเซลล์ บนชุด locked การหยุดทำแบบ offline และการตรวจตอน decode ครอบคลุมเฉพาะข้อ calibration ในงานอ่านข้อความสั้น หลายโมเดลมี precision ต่ำเพราะข้อแสดงทั้งหน้าแต่คำถามถามเพียงบางส่วน ซึ่งเรายังไม่ได้ทดสอบ

\subsection{การตีความ}
การจัดหมวดเป็นแบบอัตโนมัติ การไม่ปรากฏใน output ไม่เท่ากับโมเดลอ่านไม่ได้ และหลังการตัด การจัดหมวดแยกการวนซ้ำออกจากการอ่านไม่สำเร็จไม่ได้ เรายังไม่ได้ทดสอบบนหน้าเต็มว่าข้อความที่ไม่ปรากฏเป็นข้อความที่โมเดลอ่านได้หรือไม่
เราใช้ revision ของ Typhoon ที่เผยแพร่ภายใต้ Apache-2.0 อย่างเดียว ขณะที่เงื่อนไขปัจจุบันของผู้พัฒนาขอความยินยอมเป็นลายลักษณ์อักษรก่อนใช้เปรียบเทียบเชิงแข่งขัน

\section{สรุปผล}
ในแปดโมเดลแบบเปิดที่เราทดสอบบน benchmark เดียว mark ส่วนใหญ่ไม่ได้หายที่ตัววรรณยุกต์ บนชุด locked ทั้งแปดโมเดล mark ส่วนใหญ่ที่หายอยู่ในข้อความที่ไม่ปรากฏใน output โดยโมเดลอาจไม่พบหรือไม่ได้เขียน หรือวนซ้ำ และเราแยกไม่ได้ว่าเพราะข้อใด

เมื่อได้บริบทที่ถูกบนชุด calibration โมเดล fine-tune ภาษาไทยสองตัวเลือกวรรณยุกต์ถูกราว 99\% แม้โมเดลทั่วไปขนาด 4B ยังอ่าน mark ผิดบางตัวขณะเขียน

การหยุดเมื่อพบการวนซ้ำเป็นขั้นตอนที่ไม่ต้องฝึกเพิ่ม เพิ่มค่าประมาณ mark F1 ในทั้ง 13 เซลล์ที่ลงทะเบียน ผ่านการทดสอบที่ลงทะเบียนใน 12 เซลล์ และแทบไม่เปลี่ยน output ที่จบเอง ในการทดสอบของเรา การหักความรู้ภาษาลบ mark ที่ถูกอยู่แล้ว และไม่ได้กู้ mark ที่หายคืน ผลทั้งหมดมาจาก benchmark เดียว

\section*{หมายเหตุด้านจริยธรรมและสัญญาอนุญาต}
โมเดลทั้งหมดเป็นแบบเปิดและใช้ตามสัญญาอนุญาต (Apache-2.0 หรือ MIT โดย Typhoon ใช้ revision ที่เป็น Apache-2.0 อย่างเดียว) ThaiOCRBench เผยแพร่ภายใต้ CC-BY-SA-4.0 เราอ้างอิงบทความและ revision ของมัน รายงาน output และคะแนน และไม่เผยแพร่ภาพซ้ำ
